\section{Discussion and Conclusion}\label{sec:discussion}\label{sec:selection_method}\label{sec:conclusion}

\paragraph{When test-time scaling fails for NER.}
For discrete-answer reasoning, sample diversity reflects diverse reasoning paths; for NER, degeneracy ($12$--$68$\%) reduces candidate diversity while signal alignment collapse ($\alpha{\to}0$) limits ranking.
The degeneracy decomposition (\S\ref{sec:theory}) identifies $\alpha$, not degeneracy, as the binding constraint ($8$--$33{\times}$ tighter); this persists across families \cite{yun2025price}, scales (4B--72B), and decoding modes (Appendix~\ref{app:freeform_ablation}).
The diagnostic framework that identifies these co-occurring conditions is our primary contribution (Appendices~\ref{app:objections},~\ref{app:cost_benefit}).
At $N{=}8$, high degeneracy ($>50$\%) suggests scaling will not help; low degeneracy with strong LP-F1 alignment may warrant entity construction; otherwise, greedy remains the default (Appendices~\ref{app:recommendations},~\ref{app:when_selection_works}).
